\documentclass[12pt,a4paper]{article}
% ============================================================
% Preámbulo común a todos los documentos del Proyecto 2
% ============================================================
\usepackage[utf8]{inputenc}
\usepackage[T1]{fontenc}
\usepackage[spanish]{babel}
\usepackage[margin=2.3cm]{geometry}
\usepackage{amsmath}
\usepackage{booktabs}
\usepackage{float}
\usepackage{array}
\usepackage{enumitem}
\usepackage{tikz}
\usetikzlibrary{shapes.geometric, arrows.meta, positioning}
\usepackage[hidelinks]{hyperref}

\renewcommand{\contentsname}{Índice}
\renewcommand{\figurename}{Figura}
\renewcommand{\tablename}{Tabla}

\setlist{itemsep=2pt, topsep=4pt}

\tikzset{
  bloque/.style={rectangle, rounded corners, draw=black, thick, fill=blue!8,
                 text width=7.5cm, align=center, font=\small, inner sep=6pt},
  bloqueMCU/.style={bloque, fill=green!12},
  bloqueSPI/.style={bloque, fill=orange!25},
  bloqueSBC/.style={bloque, fill=violet!12},
  flujo/.style={-{Stealth[length=2.5mm]}, thick},
  bidir/.style={{Stealth[length=2.5mm]}-{Stealth[length=2.5mm]}, thick},
  inicio/.style={rectangle, rounded corners=6pt, draw=black, thick, fill=green!15,
                 minimum width=3.5cm, minimum height=0.8cm, align=center, font=\small},
  proc/.style={rectangle, draw=black, thick, fill=blue!8, text width=5.2cm,
               align=center, font=\small, minimum height=0.9cm, inner sep=4pt},
  dec/.style={diamond, aspect=2.4, draw=black, thick, fill=yellow!20,
              text width=3.0cm, align=center, font=\footnotesize, inner sep=1pt},
}

% Datos del proyecto (un solo lugar para mantenerlos)
\newcommand{\Universidad}{Universidad Nacional de Colombia --- Sede Manizales}
\newcommand{\Facultad}{Facultad de Ingeniería y Arquitectura}
\newcommand{\Carrera}{Carrera de Ingeniería Electrónica}
\newcommand{\Materia}{Programación de sistemas Linux embebido}
\newcommand{\TituloProyecto}{Sistema Híbrido de Adquisición Analógica mediante SPI\\[0.3cm]y Dashboard Embebido en Linux}
\newcommand{\FechaDoc}{Manizales, noviembre de 2026}

% Comando de color-neutral para valores a completar por el equipo
\newcommand{\completar}[1]{\textit{[#1]}}
% ============================================================
% Portada común: \portada{Título del documento}
% ============================================================
\newcommand{\portada}[1]{%
\begin{titlepage}
\centering
{\large \Universidad\par}
{\large \Facultad\par}
{\large \Carrera\par}
\vspace{1.2cm}
{\large Materia: \textbf{\Materia}\par}
\vspace{1.8cm}
{\LARGE\bfseries \TituloProyecto\par}
\vspace{0.5cm}
{\Large #1\par}
\vspace{1.5cm}

\small
\begin{tabular}{@{}p{0.5cm} p{4.2cm} p{2.8cm} p{6.0cm}@{}}
\toprule
& \textbf{Integrante} & \textbf{C.C.} & \textbf{Rol técnico} \\
\midrule
1. & \hrulefill & \hrulefill & Ingeniero de Documentación Técnica y Gestión de Configuración \\[0.7cm]
2. & \hrulefill & \hrulefill & Arquitecto de Sistemas Embebidos \\[0.7cm]
3. & \hrulefill & \hrulefill & Ingeniero de Software y Firmware Embebido \\[0.7cm]
4. & \hrulefill & \hrulefill & Ingeniero de Verificación y Validación (V\&V) \\[0.7cm]
\bottomrule
\end{tabular}

\vfill
\FechaDoc
\end{titlepage}
}

\begin{document}

\portada{Documento 2: Matriz de verificación y plan de pruebas}

\tableofcontents
\newpage

% ============================================================
% 1. ESTRATEGIA
% ============================================================
\section{Estrategia de diseño de experimentos}

La verificación se organiza en tres niveles. Las pruebas unitarias del firmware del ESP32-S3 validan la adquisición y la comunicación SPI de forma aislada. Las pruebas unitarias de la Raspberry Pi validan el servicio SPI y el servidor HTTP. Las pruebas de integración verifican el sistema completo, incluidos el refresco del dashboard y el rendimiento del bus a las frecuencias de muestreo requeridas.

Las pruebas de integración de tasa SPI se ejecutan a 100, 1\,000, 10\,000 y 41\,667 Sa/s. Los valores teóricos de referencia provienen del documento de arquitectura: cada ventana ocupa 320 bytes en el bus (5 segmentos de 64 bytes), es decir $t_{tx} = 2\,560 / f_{SPI}$ y ciclo de trabajo $D$ entre 0{,}02\,\% y 8{,}3\,\% a 20 MHz (0{,}04\,\% y 16{,}7\,\% a 10 MHz).

\subsection{Principio de ejecución por capas}

Se recomienda ejecutar las pruebas en este orden, y no pasar a la capa siguiente hasta que la anterior esté aprobada, porque cada una asume que la anterior funciona:

\begin{enumerate}
  \item \textbf{ESP32-S3 aislado} (solo monitor serie): presets, duración de la ventana, desfase de canales. No requiere la Raspberry Pi.
  \item \textbf{Calibración del ADC con señales conocidas:} 0 V (GND), 3{,}3 V y un divisor resistivo en cada canal para fijar ganancia y offset; después una señal periódica de frecuencia conocida (generador de funciones o un pin PWM filtrado con RC) para medir la $f_s$ real a partir de su período.
  \item \textbf{Raspberry Pi aislada:} pruebas unitarias del protocolo en el host (\texttt{make test}), configuración de \texttt{spidev}, servicio y servidor HTTP con un \texttt{window.json} de prueba.
  \item \textbf{Enlace SPI:} comenzar con \texttt{oscd --bench} a 1 MHz y subir (4, 8, 15{,}6 MHz) registrando errores de CRC y tiempos. Ajustar \texttt{--gap-us} al mínimo estable. Medir con analizador lógico en SCLK y CS.
  \item \textbf{Integración:} refresco del dashboard, cambio de preset desde la interfaz y prueba de estabilidad de 5 minutos en cada preset.
\end{enumerate}

\subsection{Instrumentación y registro}

\begin{itemize}
  \item \textbf{Analizador lógico} (por ejemplo de 8 canales a 24 MHz o más, con \texttt{sigrok}/PulseView) conectado a SCLK, CS y MISO: mide $t_{tx}$, el tiempo de guarda efectivo y la frecuencia real del reloj. Para 15{,}6 MHz se recomienda un analizador de al menos 100 MS/s; si solo se dispone de 24 MS/s se mide el período de CS y el número de flancos, no la forma del reloj.
  \item \textbf{Generador de funciones} para la verificación de $f_s$ y de ganancia. Alternativa sin instrumento: un pin PWM del propio ESP32-S3 filtrado con RC (señal de 0 a 3{,}3 V de frecuencia exacta).
  \item \textbf{Herramientas de software:} \texttt{oscd --bench N} (tiempos por transacción, ventanas perdidas, errores), \texttt{curl} para la API y el JSON, herramientas de desarrollo del navegador para los intervalos de refresco.
  \item \textbf{Registro:} cada prueba guarda su evidencia (log o captura) en \texttt{Documentos/Resultados\_Pruebas/} con el nombre del identificador de la prueba, y se actualiza la columna \emph{Estado} de la matriz con el valor medido.
\end{itemize}

\subsection{Pruebas unitarias en el ESP32-S3}

\begin{itemize}
  \item Verificación de la duración de la ventana y de la $f_s$ efectiva en cada preset.
  \item Verificación de la conversión del ADC en ambos canales y del desfase CH1/CH2.
  \item Verificación del llenado de 64 muestras por canal y del anillo de tres slots.
  \item Verificación del segmento SPI (encabezado, datos y CRC-16).
  \item Verificación de la recepción y ejecución de comandos de preset.
\end{itemize}

\subsection{Pruebas unitarias en la Raspberry Pi}

\begin{itemize}
  \item Verificación de la configuración de \texttt{spidev} (modo, velocidad y tamaño de palabra).
  \item Verificación de la validación de CRC y del descarte de tramas corruptas.
  \item Verificación de la escritura atómica del archivo de la última ventana.
  \item Verificación del servidor HTTP (Nginx) y de la entrega del JSON.
  \item Verificación del período de refresco entre 4 s y 1 s.
\end{itemize}

\subsection{Pruebas de integración}

\begin{itemize}
  \item Actualización de extremo a extremo del dashboard a 0{,}25 Hz y 1 Hz.
  \item Verificación experimental de la tasa de transferencia SPI frente a la teórica.
  \item Cambio de preset desde el dashboard.
  \item Verificación del preset de bonificación a 41\,667 Sa/s.
\end{itemize}

% ============================================================
% 2. MATRIZ
% ============================================================
\section{Matriz de verificación}

La columna de estado se completa durante la ejecución de cada prueba. Los valores teóricos provienen de la estimación del documento de arquitectura. Las pruebas \emph{Parcial} tienen una parte ejecutada (ver Sección~\ref{sec:resultados}) y otra pendiente.

\begin{table}[H]
\centering
\scriptsize
\begin{tabular}{@{}p{1.3cm} p{1.7cm} p{3.3cm} p{1.7cm} p{3.0cm} p{2.5cm} p{1.4cm}@{}}
\toprule
\textbf{ID} & \textbf{Componente} & \textbf{Descripción} & \textbf{$f_s$ (Sa/s)} & \textbf{Criterio de aceptación} & \textbf{Resultado teórico} & \textbf{Estado} \\
\midrule
UT-MCU-01 & ESP32-S3 & Duración de la ventana y $f_s$ efectiva por preset & 100; 1\,000; 10\,000; 41\,667 & $f_s$ real dentro de $\pm 1\,\%$ (con señal de frecuencia conocida) & $T_w = 64/f_s$ & Parcial \\
UT-MCU-02 & ESP32-S3 & Conversión del ADC en CH1 y CH2 y orden de canales & 1\,000 & Ganancia y offset dentro de la tolerancia del ADC; CH1 y CH2 no intercambiados & Ajuste lineal 0 a 4095 & Pendiente \\
UT-MCU-03 & ESP32-S3 & Llenado de 64 pares por ventana & 100; 1\,000; 10\,000; 41\,667 & 64 muestras por canal por ventana, sin errores del ADC & 64 pares & Aprobada \\
UT-MCU-04 & ESP32-S3 & Anillo de 3 slots bajo lectura SPI & 41\,667 & Ninguna ventana alterada durante su entrega & Cero ventanas corruptas & Aprobada \\
UT-MCU-05 & ESP32-S3 & Segmento SPI (cabecera, datos y CRC) & 1\,000 & CRC válido en 100\,\% de los segmentos & 64 bytes por segmento & Aprobada \\
UT-MCU-06 & ESP32-S3 & Ejecución de comando de preset & 100 a 41\,667 & Cambio de frecuencia en $\leq$ 1 ventana & Cambio en $\leq$ 1 $T_w$ & Aprobada \\
UT-SBC-01 & Raspberry Pi & Configuración de \texttt{spidev} & 1\,000 & Modo 0, velocidad y 8 bits confirmados & Parámetros aplicados & Aprobada \\
UT-SBC-02 & Raspberry Pi & Validación de CRC y descarte & 10\,000 & Tramas corruptas descartadas y reintentadas & 0 tramas corruptas aceptadas & Aprobada \\
UT-SBC-03 & Raspberry Pi & Escritura atómica del archivo de ventana & 1\,000 & Lectura nunca parcial desde el servidor & Archivo consistente & Aprobada \\
UT-SBC-04 & Servidor HTTP & Entrega del JSON de la última ventana & 1\,000 & HTTP 200 y JSON válido & Respuesta en $\leq$ 100 ms & Aprobada \\
UT-SBC-05 & Servicio systemd & Reinicio automático del servicio & --- & Servicio activo tras fallo simulado & Reinicio en $\leq$ 5 s & Aprobada \\
IT-01 & Sistema & Refresco del dashboard a 0{,}25 Hz & 1\,000 & Actualización cada 4 s $\pm$ 0{,}2 s & $T_r = 4$ s & Parcial \\
IT-02 & Sistema & Refresco del dashboard a 1 Hz & 1\,000 & Actualización cada 1 s $\pm$ 0{,}1 s & $T_r = 1$ s & Parcial \\
IT-03 & Enlace SPI & Transferencia de ventana a 100 Sa/s & 100 & $t_{tx}$ medido $\approx$ teórico y 0 errores de CRC & $D = 0{,}02\,\%$ (20 MHz) & Parcial \\
IT-04 & Enlace SPI & Transferencia de ventana a 1\,000 Sa/s & 1\,000 & $t_{tx}$ medido $\approx$ teórico y 0 errores de CRC & $D = 0{,}20\,\%$ (20 MHz) & Parcial \\
IT-05 & Enlace SPI & Transferencia de ventana a 10\,000 Sa/s & 10\,000 & $t_{tx}$ medido $\approx$ teórico y 0 errores de CRC & $D = 2{,}0\,\%$ (20 MHz) & Parcial \\
IT-06 & Enlace SPI & Transferencia de ventana a 41\,667 Sa/s & 41\,667 & $t_{tx}$ medido $\approx$ teórico y $D < 15\,\%$ & $D = 8{,}3\,\%$ (20 MHz) & Parcial \\
IT-07 & Sistema & Preset de bonificación de extremo a extremo & 41\,667 & Dashboard muestra la señal durante 5 min sin errores de CRC; huecos de WSEQ medidos y documentados & $f_s \geq$ 40\,000 & Aprobada \\
IT-08 & Sistema & Cambio de preset desde el dashboard & 100 a 41\,667 & La ventana siguiente al comando usa el nuevo preset en $\leq$ 2 ventanas + 1 refresco & Cambio en $\leq$ 2 $T_w$ & Parcial \\
\bottomrule
\end{tabular}
\caption{Matriz de verificación del sistema.}
\label{tab:matriz}
\end{table}

% ============================================================
% 3. UNITARIAS
% ============================================================
\section{Casos de prueba unitarios}

\subsection{ESP32-S3: adquisición y comunicación SPI}

\begin{table}[H]
\centering
\small
\begin{tabular}{@{}p{2.2cm} p{5.3cm} p{6.3cm}@{}}
\toprule
\textbf{ID} & \textbf{Procedimiento} & \textbf{Resultado esperado} \\
\midrule
UT-MCU-01 & Seleccionar cada preset (teclas 0 a 3 del monitor serie) y leer \texttt{t\_adc\_read}. Después, aplicar una señal periódica conocida y calcular $f_s$ a partir de su período en la ventana. & Duración compatible con $64/f_s$ más el arranque fijo de la ráfaga; $f_s$ real dentro de $\pm 1\,\%$. \\
UT-MCU-02 & Conectar cada canal a GND, a 3{,}3 V y a un divisor resistivo; registrar los códigos. Aplicar una tensión distinta en cada canal para comprobar el orden CH1/CH2. & Códigos acordes con el ajuste lineal; el canal de cada GPIO corresponde al índice esperado. \\
UT-MCU-03 & Contar los valores por canal en cada línea \texttt{WIN} del monitor serie. & Exactamente 64 muestras por canal en cada ventana y \texttt{err\_adc = 0}. \\
UT-MCU-04 & Con el preset de 41\,667 Sa/s, leer ventanas por SPI durante 60 s y comprobar que cada ventana recibida es consistente (misma \texttt{WSEQ} en los 5 segmentos, rampa de prueba sin discontinuidades). & Ninguna ventana mezclada o alterada. \\
UT-MCU-05 & Comparar el CRC calculado en el chip con el vector estándar y validar los segmentos recibidos en la Raspberry Pi. & CRC \texttt{0x29B1} para el vector de prueba; 100\,\% de segmentos válidos. \\
UT-MCU-06 & Enviar SET\_PRESET por SPI y medir en cuántas ventanas cambia la \texttt{fs} del encabezado. & Cambio en la ventana siguiente y sin reinicio del firmware. \\
\bottomrule
\end{tabular}
\caption{Casos de prueba unitarios del ESP32-S3.}
\end{table}

\subsection{Raspberry Pi: servicio SPI y servidor HTTP}

\begin{table}[H]
\centering
\small
\begin{tabular}{@{}p{2.2cm} p{5.3cm} p{6.3cm}@{}}
\toprule
\textbf{ID} & \textbf{Procedimiento} & \textbf{Resultado esperado} \\
\midrule
UT-SBC-01 & Ejecutar \texttt{oscd --bench 1} y leer la línea \texttt{spidev} que imprime el servicio. & Modo 0, 8 bits y velocidad máxima igual a la solicitada. \\
UT-SBC-02 & Ejecutar \texttt{make test} (host) y, en la Raspberry Pi, provocar errores en el bus (cables largos o reloj excesivo) y revisar los contadores. & Tramas corruptas descartadas y contadas; ninguna trama corrupta publicada. \\
UT-SBC-03 & Leer \texttt{window.json} en un bucle mientras el servicio lo reescribe, validando cada lectura con \texttt{python3 -m json.tool}. & Ninguna lectura parcial. \\
UT-SBC-04 & Solicitar \texttt{/data/window.json} con \texttt{curl} durante 60 s. & HTTP 200, JSON válido y respuesta en $\leq$ 100 ms. \\
UT-SBC-05 & Terminar el proceso con \texttt{kill -9}. & Reinicio automático en $\leq$ 5 s. \\
\bottomrule
\end{tabular}
\caption{Casos de prueba unitarios de la Raspberry Pi.}
\end{table}

% ============================================================
% 4. INTEGRACIÓN
% ============================================================
\section{Casos de prueba de integración}

\subsection{Refresco del dashboard (0{,}25 Hz y 1 Hz)}

\begin{enumerate}
  \item Configurar la tasa de refresco en 0{,}25 Hz desde el dashboard (o con \texttt{curl -X POST '/api/refresh?hz=0.25'}).
  \item Registrar con las herramientas de desarrollo del navegador las marcas de tiempo de cada solicitud a \texttt{window.json} durante 5 minutos.
  \item Calcular el intervalo entre solicitudes y comparar con $T_r = 4$ s.
  \item Repetir los pasos 1 a 3 con 1 Hz ($T_r = 1$ s).
  \item Criterio de aceptación: intervalo medido dentro de $\pm 0{,}2$ s (0{,}25 Hz) y $\pm 0{,}1$ s (1 Hz).
\end{enumerate}

\subsection{Verificación experimental de la tasa SPI}

\begin{enumerate}
  \item Fijar la velocidad solicitada del SPI en la Raspberry Pi y registrar la frecuencia de reloj real con el analizador lógico.
  \item Para cada preset (100, 1\,000, 10\,000 y 41\,667 Sa/s), ejecutar \texttt{oscd --bench 1000} y capturar con el analizador lógico SCLK y CS.
  \item Medir el tiempo entre el flanco de bajada y el de subida de CS para cada segmento y el tiempo entre segmentos; el de ventana es la suma de los cinco segmentos y sus guardas.
  \item Calcular el promedio y la desviación estándar de $t_{tx}$ (solo bits de reloj: $5 \times 512 / f_{SPI}$).
  \item Comparar contra $t_{tx}^{teo} = 2\,560/f_{SPI}$ y calcular el error relativo:
\end{enumerate}

\begin{equation}
\varepsilon = \frac{\bar{t}_{tx}^{exp} - t_{tx}^{teo}}{t_{tx}^{teo}} \cdot 100\,\%
\end{equation}

Adicionalmente se registra, con \texttt{oscd --bench}, el tiempo medio y máximo del \texttt{ioctl} por transacción, el número de ventanas por segundo y los contadores de ventanas perdidas, duplicadas y errores de CRC.

\subsection{Protocolo de verificación a 41\,667 Sa/s (bonificación)}

\begin{enumerate}
  \item Configurar el preset 3 en el ESP32-S3 mediante comando SPI (dashboard o \texttt{oscd --preset 3}).
  \item Comprobar con una señal de frecuencia conocida que la $f_s$ real es 41\,667 Sa/s $\pm 1\,\%$ (entonces $T_w = 1{,}536$ ms más el arranque de la ráfaga de $\approx$ 75 $\mu$s).
  \item Registrar 5 minutos de adquisición con el dashboard activo a 1 Hz.
  \item Registrar los contadores del servicio: ventanas recibidas, perdidas (saltos de \texttt{WSEQ}), errores de CRC y de sincronía. La captura es por refresco de ventanas, por lo que se espera que la Raspberry Pi no lea todas las ventanas; el criterio es \emph{cero ventanas corruptas} y que la ventana mostrada tenga la $f_s$ correcta.
  \item Verificar que el ciclo de trabajo medido de los bits es $D < 15\,\%$.
  \item Documentar cualquier desviación y su causa probable.
\end{enumerate}

% ============================================================
% 5. RESULTADOS
% ============================================================
\section{Registro de resultados}
\label{sec:resultados}

\subsection{UT-MCU-01: duración de la ventana por preset (2 de noviembre de 2026)}

Medido por software con el temporizador de ciclos del ESP32-S3 alrededor de la lectura de la ventana de 64 pares (\texttt{t\_adc\_read}), con las entradas sin señal conocida. Evidencia: \texttt{UT-MCU-01\_serial\_presets.log}.

\begin{table}[H]
\centering
\small
\begin{tabular}{@{}p{2.4cm} p{2.6cm} p{2.6cm} p{2.4cm} p{3.4cm}@{}}
\toprule
\textbf{Preset ($f_s$)} & \textbf{$T_w$ teórico} & \textbf{$T_w$ medido} & \textbf{Diferencia} & \textbf{Observación} \\
\midrule
100 Sa/s     & 640{,}000 ms & 640{,}093 ms & +0{,}09 ms (+0{,}01\,\%) & Temporizador absoluto \\
1\,000 Sa/s  & 64{,}000 ms  & 64{,}318 ms  & +0{,}32 ms (+0{,}50\,\%) & Ráfaga DMA \\
10\,000 Sa/s & 6{,}400 ms   & 6{,}492 ms   & +0{,}09 ms (+1{,}4\,\%)  & Ráfaga DMA \\
41\,667 Sa/s & 1{,}536 ms   & 1{,}609 ms   & +0{,}07 ms (+4{,}8\,\%)  & Ráfaga DMA \\
\bottomrule
\end{tabular}
\caption{Duración medida de la ventana. La diferencia en los presets DMA es compatible con un tiempo fijo de arranque de la ráfaga de 75 a 90 $\mu$s.}
\end{table}

La prueba queda \textbf{parcial}: confirma que la ráfaga completa de 64 pares dura lo esperado, pero la $f_s$ efectiva por canal debe confirmarse con una señal de frecuencia conocida (ver procedimiento). Con este resultado se confirma que el ADC del ESP32-S3 alcanza el preset de 41\,667 Sa/s por canal (83\,333 conversiones por segundo en total).

\subsection{UT-MCU-03, UT-MCU-05, UT-MCU-06 y UT-SBC-02 (parciales)}

\begin{itemize}
  \item \textbf{UT-MCU-03 (aprobada):} cada línea \texttt{WIN} del monitor serie contiene 64 valores por canal y el contador \texttt{err\_adc} permaneció en 0 durante todas las capturas.
  \item \textbf{UT-MCU-05 (parcial):} el CRC-16 del vector \texttt{"123456789"} calculado en el chip es \texttt{0x29B1} (evidencia \texttt{UT-MCU-05\_crc\_autotest.log}). Falta validar segmentos reales recibidos por la Raspberry Pi.
  \item \textbf{UT-MCU-06 (parcial):} el cambio de preset por teclado del monitor serie se refleja en la ventana siguiente en los cuatro presets. Falta repetir la prueba con el comando SPI.
  \item \textbf{UT-SBC-02 (parcial):} el protocolo de la Raspberry Pi pasa 11 comprobaciones en el host, entre ellas el descarte de una trama con un bit alterado y una con sincronía errónea (evidencia \texttt{UT-SBC-02\_test\_protocol\_host.log}). Falta la prueba sobre el bus real.
\end{itemize}

\subsection{Pruebas en la Raspberry Pi 4 sin el ESP32-S3 (2 de noviembre de 2026)}

Ejecutadas en una Raspberry Pi 4 Model B Rev 1.2 (kernel 6.18). Las pruebas de servidor, API y servicio usan el modo \texttt{oscd --demo}, que genera ventanas sintéticas con el mismo formato JSON; no prueban el enlace SPI. Evidencia: \texttt{UT-SBC-01\_spidev\_ioctl.log} y \texttt{UT-SBC-03-05\_IT-01-02\_raspberry.log}.

\begin{table}[H]
\centering
\small
\begin{tabular}{@{}p{1.8cm} p{6.4cm} p{6.8cm}@{}}
\toprule
\textbf{ID} & \textbf{Resultado medido} & \textbf{Conclusión} \\
\midrule
UT-SBC-01 & \texttt{spidev0.0} acepta modo 0, 8 bits y las velocidades pedidas (1, 8, 15{,}6, 20 y 31{,}25 MHz) & Aprobada. El reloj efectivo queda por medir con analizador lógico \\
UT-SBC-03 & 189\,609 lecturas directas de \texttt{window.json} mientras el servicio lo reescribía: 0 parciales, 25 versiones distintas vistas & Aprobada: la escritura con \texttt{rename} es atómica \\
UT-SBC-04 & 300 solicitudes HTTP (una cada 100 ms): 300 con código 200 y JSON válido; tiempo máximo 48{,}7 ms & Aprobada ($\leq$ 100 ms) \\
UT-SBC-05 & \texttt{kill -9} al servicio: nuevo PID y respuesta HTTP 200 en 1{,}09 s & Aprobada ($\leq$ 5 s) \\
IT-01 & Escritura del JSON a 0{,}25 Hz: 8 intervalos, media 4{,}000 s, rango 4{,}000 a 4{,}001 s & Parcial: falta medir las solicitudes del navegador \\
IT-02 & Escritura del JSON a 1 Hz: 11 intervalos, media 1{,}001 s, rango 1{,}000 a 1{,}001 s & Parcial: falta medir las solicitudes del navegador \\
\bottomrule
\end{tabular}
\caption{Pruebas ejecutadas en la Raspberry Pi 4.}
\end{table}

Tiempo del \texttt{ioctl} de una transacción de 64 bytes (sin ESP32-S3, MISO flotante), promedio de miles de transacciones:

\begin{table}[H]
\centering
\small
\begin{tabular}{@{}p{3.4cm} p{3.4cm} p{3.4cm} p{3.4cm}@{}}
\toprule
\textbf{Velocidad solicitada} & \textbf{Tiempo ideal de bits} & \textbf{\texttt{ioctl} medio} & \textbf{\texttt{ioctl} máximo} \\
\midrule
1 MHz      & 512{,}0 $\mu$s & 590{,}5 $\mu$s & 625{,}3 $\mu$s \\
8 MHz      & 64{,}0 $\mu$s  & 96{,}5 $\mu$s  & 146{,}9 $\mu$s \\
15{,}625 MHz & 32{,}8 $\mu$s & 58{,}7 $\mu$s & 101{,}7 $\mu$s \\
20 MHz     & 25{,}6 $\mu$s  & 52{,}6 $\mu$s  & 111{,}8 $\mu$s \\
31{,}25 MHz & 16{,}4 $\mu$s & 29{,}5 $\mu$s  & 104{,}2 $\mu$s \\
\bottomrule
\end{tabular}
\caption{Latencia del \texttt{ioctl} SPI en la Raspberry Pi 4.}
\end{table}

El \texttt{ioctl} añade entre 13 y 78 $\mu$s sobre el tiempo ideal de bits. Con cinco transacciones por ventana y 150 $\mu$s de guarda, un ciclo de lectura cuesta aproximadamente 1{,}0 ms a 20 MHz y 1{,}2 ms a 8 MHz, frente a una ventana nueva cada $\approx$ 1{,}6 ms en el preset de 41\,667 Sa/s: el margen existe pero es estrecho y la guarda mínima estable debe determinarse con el ESP32-S3 conectado.

\subsection{Integración ESP32-S3 + Raspberry Pi 4 (5 de noviembre de 2026)}

Montaje: ESP32-S3-DevKitC-1 (WROOM-1 N8R2) cableado a la Raspberry Pi 4 según el Documento 1, cables cortos y GND común. Medición con \texttt{oscd --bench}, que cuenta ventanas recibidas, ventanas no leídas (saltos de WSEQ), duplicadas, errores de CRC y de sincronía. Evidencia: \texttt{IT-03-08\_integracion\_spi.log}.

\paragraph{Velocidad del reloj SPI} (preset de 41\,667 Sa/s, guarda de 150 $\mu$s, 3\,000 ventanas o 8 s):

\begin{table}[H]
\centering
\small
\begin{tabular}{@{}p{2.6cm} p{2.4cm} p{2.6cm} p{2.2cm} p{2.0cm} p{2.2cm}@{}}
\toprule
\textbf{Reloj solicitado} & \textbf{Ventanas/s} & \textbf{\texttt{ioctl} medio} & \textbf{No leídas} & \textbf{Errores CRC} & \textbf{Errores sync} \\
\midrule
1 MHz        & 184{,}7 & 809{,}6 $\mu$s & 3\,010 & 0 & 337 \\
8 MHz        & 579{,}8 & 165{,}1 $\mu$s & 112    & 0 & 553 \\
15{,}625 MHz & 594{,}2 & 91{,}8 $\mu$s  & 37     & 0 & 612 \\
20 MHz       & 585{,}9 & 79{,}0 $\mu$s  & 6      & 0 & 883 \\
\bottomrule
\end{tabular}
\caption{Desempeño del enlace según el reloj solicitado.}
\end{table}

\paragraph{Tiempo de guarda} (20 MHz, preset de 41\,667 Sa/s): con 50 $\mu$s el esclavo no alcanza a rearmarse y se pierden muchas ventanas; con 100 y 150 $\mu$s se reciben 587 a 590 ventanas/s con menos de 1{,}1\,\% no leídas; con 250 $\mu$s baja a 554 ventanas/s y con 400 $\mu$s a 382 ventanas/s. Se adoptó \textbf{20 MHz y 150 $\mu$s} como configuración del servicio. Las corridas de velocidad y de guarda se hicieron antes de corregir los defectos de la Sección~\ref{sec:defectos}, por lo que algunos conteos de ventanas no leídas (en especial el de 50 $\mu$s) pueden estar inflados; los errores de CRC no se ven afectados.

\paragraph{Presets} (20 MHz, 150 $\mu$s; el cambio se ordena por SPI al inicio de cada corrida):

\begin{table}[H]
\centering
\small
\begin{tabular}{@{}p{2.4cm} p{2.6cm} p{2.6cm} p{3.4cm} p{2.6cm}@{}}
\toprule
\textbf{Preset} & \textbf{Ventanas/s recibidas} & \textbf{Ventanas/s esperadas} & \textbf{Tiempo hasta ver el preset nuevo} & \textbf{Errores CRC} \\
\midrule
100 Sa/s      & 1{,}6   & 1{,}56 & 649 ms & 0 \\
1\,000 Sa/s   & 15{,}0  & 15{,}5 & 661 ms & 0 \\
10\,000 Sa/s  & 146{,}1 & 154    & 29 ms  & 0 \\
41\,667 Sa/s  & 587{,}0 & 621    & 11 ms  & 0 \\
\bottomrule
\end{tabular}
\caption{Desempeño por preset. Las ventanas/s esperadas son $1/T_w$ con la duración medida en UT-MCU-01 (incluye el arranque de la ráfaga); entre ráfagas el firmware añade además el tiempo de publicar la ventana, por lo que a 41\,667 Sa/s el ESP32-S3 produce en la práctica $\approx$ 595 ventanas/s.}
\end{table}

El tiempo de cambio incluye terminar la ventana en curso del preset anterior: al pasar de 100 Sa/s a 1\,000 Sa/s debe terminar una ventana de 640 ms, de ahí los 661 ms. En todos los casos el cambio ocurre en la ventana siguiente a la que estaba en curso (UT-MCU-06).

\paragraph{IT-07: estabilidad de 5 minutos a 41\,667 Sa/s} con el servicio \texttt{oscd} (20 MHz, 150 $\mu$s) y lectura de \texttt{window.json} por Nginx cada segundo:

\begin{itemize}
  \item Duración: 332 s; 194\,668 ventanas recibidas (586{,}3 ventanas/s).
  \item \textbf{Errores de CRC: 0}; ventanas hacia atrás: 0; errores de E/S: 0; reinicios del servicio: 0.
  \item Fracción de ventanas producidas que se leyeron: \textbf{99{,}6\,\%} (783 no leídas).
  \item Lecturas HTTP del JSON: 300 de 300 correctas, todas con preset 3 y 64 muestras por canal.
  \item Errores de sincronía: 54\,274 (segmentos clocados cuando el esclavo aún no estaba rearmado); se descartan sin consecuencias y se recupera en el siguiente segmento 0.
\end{itemize}

\paragraph{IT-03 a IT-06 (parciales).} El enlace transfiere ventanas en los cuatro presets con 0 errores de CRC. El tiempo medio del \texttt{ioctl} a 20 MHz es de $\approx$ 79 $\mu$s por segmento ($\approx$ 395 $\mu$s por ventana), que incluye la sobrecarga del kernel. Falta medir con analizador lógico la frecuencia real del reloj y $t_{tx}$ de los bits, para comparar con el valor teórico $2\,560/f_{SPI} = 128\ \mu$s.

\paragraph{IT-08 (parcial).} El cambio de preset a través de la API publicada por Nginx (\texttt{POST /api/preset}, la misma llamada que usan los botones del dashboard) se reflejó en \texttt{window.json} en menos de 2 s. Falta comprobarlo pulsando los botones en un navegador.

\subsection{Defectos encontrados y corregidos durante la integración}
\label{sec:defectos}

\begin{table}[H]
\centering
\small
\begin{tabular}{@{}p{3.6cm} p{5.6cm} p{5.6cm}@{}}
\toprule
\textbf{Defecto} & \textbf{Causa} & \textbf{Corrección} \\
\midrule
El preset no cambiaba al iniciar \texttt{oscd} & El número del comando se asignaba antes de conocer el \texttt{CMD\_ACK} del esclavo; si coincidía con un comando ya ejecutado, el ESP32-S3 lo ignoraba y la Raspberry Pi lo daba por confirmado & El número se asigna después de alinearse con \texttt{CMD\_ACK} y nunca se envía antes \\
CH1 siempre en 0 a 100 Sa/s y primera muestra de CH1 en 0 en las ráfagas & La primera conversión de cada ráfaga DMA del ADC vale 0 & Se adquiere un par adicional al inicio de cada ráfaga y se descarta \\
Salto falso de miles de ventanas «no leídas» al arrancar el servicio & Sin lecturas, el esclavo deja congelada la ventana de ese momento, que puede tener decenas de segundos & Se descarta la primera ventana tras arrancar; los saltos hacia atrás se cuentan aparte y no se publican \\
\bottomrule
\end{tabular}
\caption{Defectos de integración.}
\end{table}

Tras las correcciones, un reinicio del servicio después de 20 s sin lecturas dio 0 ventanas no leídas.

\subsection{Hallazgo abierto sobre las entradas}

Con las entradas sin una señal conocida, ambos canales muestran a 100 Sa/s un patrón periódico de 5 muestras entre 0 y 4095, compatible con entradas flotantes que captan la red eléctrica (alias de 20 Hz de un armónico de 120 Hz). UT-MCU-02 sigue pendiente: conectar GND, 3{,}3 V y un divisor a cada canal, con una tensión distinta en cada uno para confirmar el orden CH1/CH2.

\end{document}
